% GED Math Unit: Equations, Inequalities, and Functions
% Beamer Slide Deck

\documentclass{beamer}
\usetheme{Madrid}
\usecolortheme{seahorse}
\usepackage{amsmath, amssymb}
\title{GED Math: Equations, Inequalities, and Functions}

\date{\today}

\begin{document}

% Title Slide
\begin{frame}
  \titlepage
\end{frame}

%=============================
% LESSON 1: LINEAR EQUATIONS
%=============================
\begin{frame}{Lesson 1: Solving Linear Equations}
\begin{itemize}
    \item Understand the structure of linear equations.
    \item Apply inverse operations to isolate variables.
    \item Check your solution by substitution.
\end{itemize}
\end{frame}

\begin{frame}{Worked Examples}
\begin{enumerate}
  \item $2x + 5 = 13 \Rightarrow 2x = 8 \Rightarrow x = 4$
  \item $5x - 3 = 2x + 6 \Rightarrow 3x = 9 \Rightarrow x = 3$
  \item $7 - 2x = 11 \Rightarrow -2x = 4 \Rightarrow x = -2$
  \item $3(x + 2) = 15 \Rightarrow x + 2 = 5 \Rightarrow x = 3$
  \item $\frac{x}{4} + 2 = 5 \Rightarrow \frac{x}{4} = 3 \Rightarrow x = 12$
\end{enumerate}
\end{frame}

%=============================
% LESSON 2: LINEAR INEQUALITIES
%=============================
\begin{frame}{Lesson 2: Solving Inequalities}
\begin{itemize}
  \item Solve inequalities similarly to equations.
  \item Flip the inequality when multiplying/dividing by a negative number.
  \item Represent solutions on a number line.
\end{itemize}
\end{frame}

\begin{frame}{Worked Examples}
\begin{enumerate}
  \item $x + 3 < 7 \Rightarrow x < 4$
  \item $2x - 5 \geq 1 \Rightarrow 2x \geq 6 \Rightarrow x \geq 3$
  \item $-3x \leq 9 \Rightarrow x \geq -3$ (flip inequality)
  \item $5 - x > 2 \Rightarrow -x > -3 \Rightarrow x < 3$
  \item $\frac{x}{2} + 4 \leq 6 \Rightarrow \frac{x}{2} \leq 2 \Rightarrow x \leq 4$
\end{enumerate}
\end{frame}

%=============================
% LESSON 3: SYSTEMS OF EQUATIONS
%=============================
\begin{frame}{Lesson 3: Systems of Linear Equations}
\begin{itemize}
  \item Solve by substitution or elimination.
  \item Interpret intersection point as common solution.
\end{itemize}
\end{frame}

\begin{frame}{Worked Examples}
\begin{enumerate}
  \item \( x + y = 10, \ y = 4 \Rightarrow x = 6 \)
  \item \( 2x + y = 7, \ x - y = 1 \Rightarrow (x, y) = (\tfrac{8}{3}, \tfrac{5}{3}) \)
  \item \( 3x - y = 5, \ 2x + y = 4 \Rightarrow (x, y) = (\tfrac{9}{5}, -\tfrac{2}{5}) \)
  \item \( y = 2x + 1, \ y = -x + 4 \Rightarrow (x, y) = (1, 3) \)
  \item \( x + 2y = 8, \ 3x - y = 5 \Rightarrow (x, y) = (3, 2.5) \)
\end{enumerate}
\end{frame}

%=============================
% LESSON 4: FUNCTIONS AND GRAPHS
%=============================
\begin{frame}{Lesson 4: Functions and Their Representations}
\begin{itemize}
  \item A function assigns exactly one output to each input.
  \item Use function notation: $f(x) = 2x + 3$.
  \item Identify domain and range.
\end{itemize}
\end{frame}

\begin{frame}{Worked Examples}
\begin{enumerate}
  \item $f(x) = 2x + 3, \ f(4) = 11$
  \item $g(x) = -x^2 + 4, \ g(2) = 0$
  \item $h(x) = 3x - 5, \ h(-1) = -8$
  \item $f(x) = |x - 2|, \ f(5) = 3$
  \item Determine if $\{(1,2), (2,3), (3,4), (3,5)\}$ is a function — No (3 maps to two values)
\end{enumerate}
\end{frame}

%=============================
% LESSON 5: QUADRATIC FUNCTIONS
%=============================
\begin{frame}{Lesson 5: Quadratic Functions}
\begin{itemize}
  \item Standard form: $y = ax^2 + bx + c$
  \item Axis of symmetry: $x = -\frac{b}{2a}$
  \item Vertex: $(h, k)$
\end{itemize}
\end{frame}

\begin{frame}{Worked Examples}
\begin{enumerate}
  \item $y = x^2 - 4x + 3 \Rightarrow$ vertex at $(2, -1)$
  \item $y = -x^2 + 2x + 3 \Rightarrow$ vertex at $(1, 4)$
  \item $y = 2(x - 1)^2 + 3$ vertex form $(1, 3)$
  \item Solve $x^2 - 5x + 6 = 0 \Rightarrow (x - 2)(x - 3) = 0 \Rightarrow x = 2,3$
  \item Axis of symmetry of $y = x^2 + 6x + 5$ is $x = -3$
\end{enumerate}
\end{frame}

\end{document}
